\section{Журнал выполнения}
\begin{enumerate}
    \item Что сжимать. В индексе ЛР7 на корпусе в 1 млн документов 88\,\% файла --- списки документов, поэтому сжимались только они. Словарь (таблица термов 61 МиБ и строки термов) и прямой индекс остались как были.
    \item Порядок байтов в VByte. В лекции описан вариант, где бит $c = 1$ стоит в последнем байте числа, а группы по 7 бит идут начиная с младших. Так и сделано: при распаковке группа сразу сдвигается на своё место, длину числа заранее знать не нужно.
    \item Запись файла. Таблица термов с полем \texttt{first} записана в файле раньше списков, а при сжатии \texttt{first} --- смещение в байтах, которое зависит от длины кодов всех предыдущих списков. Поэтому при записи два прохода по спискам: сначала считаются длины кодов (без записи), затем списки кодируются прямо в буфер записи. Второй вариант --- кодировать в память и потом писать --- потребовал бы ещё сотни мегабайт памяти.
    \item Защита от повреждённых данных. Несжатый список нельзя прочитать \enquote{неправильно}: каждый docID --- ровно 4 байта. У сжатого испорченный байт сдвигает все следующие числа. Поэтому распаковка проверяет, что список кончается ровно на своей границе, числа не длиннее 5 байт и не выходят за 32 бита, разности положительны; при загрузке проверяется, что смещения \texttt{first} не убывают. Читатель теперь отказывается открывать индекс с неизвестным флагом: иначе старая программа прочитала бы сжатый индекс как несжатый.
    \item Совместимость. Ключ \texttt{-c} необязателен, индекс ЛР7 без сжатия (и индексы ЛР4) читаются как раньше; автотесты ЛР4, ЛР7 и ЛР8 после изменений пройдены (10, 8 и 7 тестов).
    \item Замеры. Процессор гибридный (P-, E- и LP E-ядра), время одних и тех же действий меняется до 2 раз, поэтому везде берётся минимум из нескольких запусков. Индексация 1 млн документов идёт около 9 минут, больше половины времени --- чтение 1 млн файлов, и оно зависит от того, сколько файлов осталось в кэше ОС после предыдущего запуска. Сжатие меняет только этап записи индекса, поэтому время сравнивается по этапам.
    \item Проверка на 30 тыс. документов (корпус ЛР4) перед полным прогоном: списки 52 МиБ $\to$ 16 МиБ, выдача 24 запросов ЛР8 совпала побайтно.
\end{enumerate}

\section{План тестирования}
Корректность поиска после сжатия следует из двух утверждений: (1) кодек обратим --- распаковка кода VByte любого возрастающего списка docID возвращает тот же список; (2) индекс со сжатием содержит те же термы, $df$ и списки, что и без сжатия. Поиск ЛР8 получает от \texttt{read\_list} тот же массив docID, что и раньше, а операции над массивами не менялись, поэтому выдача та же. Тесты проверяют оба утверждения по отдельности и затем результат целиком --- выдачу поиска.

Юнит-тесты на C++ (\texttt{tests/test\_compress.cpp}) на каркасе \texttt{lab3/tests/unittest.h}; запуск --- \texttt{make test}. Эталоны в тестах считаются независимо от проверяемого кода (контейнеры STL, свой разбор файла по побитовой схеме).

\begin{center}
\small
\begin{tabular}{|c|p{0.3\textwidth}|p{0.58\textwidth}|}
\hline
\textnumero & Тест & Что проверяется \\
\hline
1 & Схема одного числа & Байты кода для 0, 5, 127, 128, 130, 824, 16\,383, 16\,384, 214\,577, $2^{32} - 1$ совпадают с выписанными вручную по схеме. \\
\hline
2 & Длина и обратимость & На границах $2^7$, $2^{14}$, $2^{21}$, $2^{28}$ длина кода 1--5 байт; $10^6$ случайных чисел всех длин: распаковка возвращает число, бит конца стоит только в последнем байте. \\
\hline
3 & Список документов & Пример лекции 824, 829, 215\,406 $\to$ \texttt{38 86 85 31 0c 8d}; docID 0 и сплошной список; 2000 случайных списков с разностями от 1 до 3 млн и docID до $2^{32} - 1$ упаковываются и распаковываются без потерь. \\
\hline
4 & Повреждённые данные & Распаковка отвергает: нет байта конца, число длиннее 5 байт или больше $2^{32} - 1$, нулевую разность, переполнение docID, лишние и недостающие байты. \\
\hline
5 & Файл по побитовой схеме & Индекс маленького корпуса, построенный с \texttt{-c}, разбирается в тесте своим кодом по описанию формата: флаг 2, смещения \texttt{first} в байтах, списки подряд, 11 записей --- 11 байт; остальные разделы совпадают с несжатым индексом байт в байт. \\
\hline
6 & \texttt{read\_list} на случайном корпусе & 50\,000 документов: все списки сжатого и несжатого индексов совпадают, включая терм во всех документах (50\,000 байт) и разности 20\,000 и 49\,999 (по 3 байта). \\
\hline
7 & 500 случайных запросов & Случайные запросы с И, ИЛИ, НЕ и скобками по индексу на 20\,000 документов: выдача по сжатому индексу совпадает с выдачей по несжатому (а та в ЛР8 сверена с вычислением на множествах). \\
\hline
8 & Стемминг и сжатие & Флаги 1 и 2 вместе (\texttt{flags} = 3): поиск со стеммингом по сжатому индексу. \\
\hline
9 & Повреждённый индекс & Снят бит конца в списке --- этот список не читается, остальные читаются, \texttt{search} сообщает о повреждении; неизвестный флаг и убывающие смещения --- индекс не открывается. \\
\hline
10 & Утилиты & \texttt{indexer -c}, \texttt{lookup}, \texttt{search} и \texttt{listbench} как отдельные процессы: одинаковые списки и выдача по обоим индексам, \texttt{listbench -v} --- 0 расхождений. \\
\hline
\end{tabular}
\end{center}

Проверки на индексе всего корпуса (\texttt{scripts/bench.py}, \texttt{results/verify.txt}):
\begin{itemize}
    \item \texttt{listbench -v}: у сжатого индекса те же термы, $df$, списки (все 364 млн записей) и записи документов, что у индекса ЛР7;
    \item сумма docID каждой группы термов при замере скорости одинакова для обоих индексов;
    \item 24 запроса ЛР8 с выдачей всех результатов (\texttt{-n 0}): выдача по обоим индексам совпадает построчно (сравнивается SHA-256).
\end{itemize}

\section{Результаты}
\subsection{Автотесты}
Пройдены все 10 тестов. Журнал прогона (\texttt{results/tests.log}):

{\scriptsize
\VerbatimInput[breaklines, breakanywhere]{../results/tests.log}
}

Проверки на индексе корпуса (\texttt{results/verify.txt}):

{\scriptsize
\VerbatimInput[breaklines, breakanywhere]{../results/verify.txt}
}

\subsection{Размер индекса}
Размеры разделов индекса корпуса (1\,000\,000 документов, 3\,974\,421 терм, 363\,958\,293 записи):
\begin{center}
\small
\begin{tabular}{|l|r|r|r|}
\hline
Раздел & Без сжатия, байт & VByte, байт & VByte / без сжатия, \% \\
\hline
Заголовок & 64 & 64 & 100,0 \\
\hline
Таблица документов & 12\,000\,000 & 12\,000\,000 & 100,0 \\
\hline
Таблица термов & 63\,590\,736 & 63\,590\,736 & 100,0 \\
\hline
Списки документов & 1\,455\,833\,172 & 439\,604\,065 & 30,2 \\
\hline
Пул строк & 129\,427\,722 & 129\,427\,722 & 100,0 \\
\hline
Всего & 1\,660\,851\,694 & 644\,622\,587 & 38,8 \\
\hline
\end{tabular}

\end{center}
Списки документов уменьшились в 3,3 раза (с 1\,388 до 419 МиБ), в среднем 9,66 бита на запись вместо 32; весь индекс --- в 2,6 раза (с 1\,584 до 615 МиБ). Остальные разделы не менялись, и теперь они составляют уже треть файла: словарь со строками термов и прямой индекс стоило бы сжимать следующими.

Сколько бит на запись тратит VByte в зависимости от $df$ (u32 --- без сжатия):
\begin{center}
\footnotesize
\setlength{\tabcolsep}{4pt}
\begin{tabular}{|l|r|r|r|r|r|r|r|r|r|}
\hline
$df$ & Термов & Записей & \multicolumn{2}{c|}{МиБ} & Бит & \multicolumn{2}{c|}{мс} & \multicolumn{2}{c|}{нс/зап.} \\
\hline
 &  &  & u32 & VB & /зап. & u32 & VB & u32 & VB \\
\hline
1 & 2\,382\,732 & 2\,382\,732 & 9,1 & 6,7 & 23,61 & 15,9 & 17,9 & 6,66 & 7,51 \\
\hline
2--10 & 1\,190\,598 & 4\,414\,325 & 16,8 & 11,8 & 22,34 & 20,8 & 30,6 & 4,71 & 6,93 \\
\hline
11--100 & 308\,118 & 9\,169\,089 & 35,0 & 20,2 & 18,52 & 18,7 & 46,4 & 2,04 & 5,07 \\
\hline
101--$10^3$ & 67\,439 & 20\,737\,071 & 79,1 & 37,0 & 14,96 & 20,9 & 53,4 & 1,01 & 2,58 \\
\hline
$10^3$--$10^4$ & 19\,873 & 60\,848\,498 & 232,1 & 83,7 & 11,53 & 41,5 & 196,8 & 0,68 & 3,23 \\
\hline
$10^4$--$10^5$ & 5\,136 & 149\,965\,910 & 572,1 & 148,7 & 8,32 & 87,8 & 165,1 & 0,59 & 1,10 \\
\hline
$> 10^5$ & 525 & 116\,440\,668 & 444,2 & 111,2 & 8,01 & 59,5 & 92,3 & 0,51 & 0,79 \\
\hline
\end{tabular}

\end{center}
У термов, встретившихся один раз, docID до $10^6$ записывается почти всегда 3 байтами (23,6 бита). С ростом $df$ разности уменьшаются, и у термов с $df > 10^5$ (средняя разность меньше 10) почти каждая запись --- ровно 1 байт, 8,01 бита. Меньше 8 бит VByte не даёт, а именно здесь находятся 73\,\% всех записей.

\subsection{Скорость прохода по координатным блокам}
Проход --- чтение всех списков группы в массив docID (\texttt{read\_list}, как в поиске) и сумма docID; индекс уже в памяти. Время --- минимум из 10 проходов:
\begin{center}
\footnotesize
\setlength{\tabcolsep}{4pt}
\begin{tabular}{|p{0.25\textwidth}|r|r|r|r|r|r|r|r|r|}
\hline
Группа & Термов & Записей & \multicolumn{2}{c|}{МиБ} & Бит & \multicolumn{2}{c|}{мс} & \multicolumn{2}{c|}{нс/зап.} \\
\hline
 &  &  & u32 & VB & /зап. & u32 & VB & u32 & VB \\
\hline
Все термы & 3\,974\,421 & 363\,958\,293 & 1\,388,4 & 419,2 & 9,66 & 224,4 & 582,1 & 0,62 & 1,60 \\
\hline
Редкие, $df \le 100$ & 3\,881\,448 & 15\,966\,146 & 60,9 & 38,7 & 20,34 & 42,3 & 85,1 & 2,65 & 5,33 \\
\hline
Средние, $100 < df \le 10^4$ & 87\,312 & 81\,585\,569 & 311,2 & 120,6 & 12,40 & 58,4 & 246,1 & 0,72 & 3,02 \\
\hline
Частые, $df > 10^4$ & 5\,661 & 266\,406\,578 & 1\,016,3 & 259,9 & 8,18 & 149,1 & 257,0 & 0,56 & 0,96 \\
\hline
\end{tabular}

\end{center}
\begin{itemize}
    \item \textbf{Все термы:} проход стал в 2,6 раза дольше (224 $\to$ 582 мс, 0,62 $\to$ 1,60 нс на запись) при списках в 3,3 раза меньше.
    \item \textbf{Редкие} ($df \le 100$): сжатие хуже всего (64\,\% размера: 20,3 бита на запись), замедление в 2 раза. Время на запись большое у обоих индексов: на 3,9 млн коротких списков приходится всего 16 млн записей, и основное время --- переход к очередному списку (промах кэша), а не чтение чисел.
    \item \textbf{Средние} ($100 < df \le 10^4$): 39\,\% размера, но наибольшее замедление --- в 4,2 раза (0,72 $\to$ 3,02 нс). Средняя разность здесь от 100 до $10^4$, коды в 1 и 2 байта перемешаны, и процессор не может предсказать ветвление \enquote{последний ли это байт числа}.
    \item \textbf{Частые} ($df > 10^4$): лучшее сжатие (26\,\% размера, 8,18 бита) и наименьшее замедление --- в 1,7 раза (0,56 $\to$ 0,97 нс): почти все коды однобайтовые, ветвление предсказуемо.
\end{itemize}

\subsection{Скорость индексации}
Весь корпус индексировался 4 раза в порядке \enquote{без сжатия, VByte, VByte, без сжатия}; по каждому этапу --- минимум из двух запусков (\texttt{results/indexing.tsv}):
\begin{center}
\small
\begin{tabular}{|p{0.5\textwidth}|r|r|}
\hline
Показатель & Без сжатия & VByte \\
\hline
Размер индекса, МиБ & 1\,583,9 & 614,8 \\
\hline
\quad из них списки документов & 1\,388,4 & 419,2 \\
\hline
Время индексации (минимум), с & 151,6 & 166,4 \\
\hline
\quad чтение файлов & 79,9 & 90,4 \\
\hline
\quad токенизация и вставка в словарь & 67,5 & 71,2 \\
\hline
\quad сортировка словаря & 1,88 & 1,91 \\
\hline
\quad запись индекса (со сжатием списков) & 2,11 & 2,60 \\
\hline
Пиковая память, МиБ & 2\,184 & 2\,184 \\
\hline
\end{tabular}

\end{center}
Сжатие меняет только этап записи: он стал дольше на 0,5 с (2,1 $\to$ 2,6 с; добавились подсчёт длин кодов и кодирование 364 млн разностей), это 0,3\,\% времени индексации. Остальные этапы код не затрагивает, их разница --- шум: время всего прогона на одних и тех же данных было 284, 186, 166 и 152 с, оно зависит в основном от того, сколько из 1 млн файлов (5,6 ГиБ текста) осталось в кэше ОС. Пиковая память та же, 2,1 ГиБ: списки копятся в памяти несжатыми, сжимаются только при записи.

\subsection{Скорость поиска}
24 запроса ЛР8, каждый повторён 20 раз, минимум из 3 запусков; время запроса не включает загрузку индекса и печать выдачи:
\begin{center}
\small
\begin{tabular}{|p{0.45\textwidth}|r|r|}
\hline
Показатель & Без сжатия & VByte \\
\hline
Загрузка индекса в память, мс & 840 & 348 \\
\hline
Запрос: медиана, мкс & 645,2 & 727,2 \\
\hline
Запрос: среднее, мкс & 2\,616,8 & 2\,803,4 \\
\hline
Запрос: максимум, мкс & 23\,611,6 & 24\,114,4 \\
\hline
\end{tabular}

\end{center}
{\small
\begin{longtable}{|r|p{0.5\textwidth}|r|r|r|}
\hline
\textnumero & Запрос после разбора & Найдено & u32, мкс & VB, мкс \\
\hline
\endhead
1 & \texttt{lviv} & 1\,508 & 1,0 & 3,6 \\
\hline
2 & \texttt{museum} & 124\,133 & 34,2 & 81,5 \\
\hline
3 & \texttt{the} & 987\,676 & 277,7 & 678,6 \\
\hline
4 & \texttt{moscow \&\& aviation \&\& institute} & 475 & 143,0 & 204,3 \\
\hline
5 & \texttt{cheap \&\& hostel \&\& paris} & 496 & 239,6 & 314,3 \\
\hline
6 & \texttt{beach \&\& surfing \&\& australia} & 672 & 135,9 & 204,5 \\
\hline
7 & \texttt{the \&\& and \&\& of \&\& in \&\& a} & 905\,478 & 10\,413,6 & 10\,300,0 \\
\hline
8 & \texttt{ski \textbar{}\textbar{} snowboard} & 13\,941 & 16,4 & 40,7 \\
\hline
9 & \texttt{castle \textbar{}\textbar{} palace \textbar{}\textbar{} fortress} & 90\,555 & 515,3 & 577,1 \\
\hline
10 & \texttt{the \textbar{}\textbar{} and \textbar{}\textbar{} of \textbar{}\textbar{} in \textbar{}\textbar{} a \textbar{}\textbar{} to} & 999\,794 & 6\,954,4 & 8\,096,1 \\
\hline
11 & \texttt{museum \&\& !art} & 75\,980 & 801,5 & 882,3 \\
\hline
12 & \texttt{beach \&\& !sand \&\& !swimming} & 62\,268 & 502,5 & 580,8 \\
\hline
13 & \texttt{!lviv} & 998\,492 & 775,1 & 775,7 \\
\hline
14 & \texttt{!the} & 12\,324 & 1\,000,3 & 1\,356,2 \\
\hline
15 & \texttt{!(the \textbar{}\textbar{} and)} & 2\,539 & 5\,206,5 & 5\,244,1 \\
\hline
16 & \texttt{(red \textbar{}\textbar{} yellow) \&\& car} & 20\,971 & 920,1 & 1\,020,9 \\
\hline
17 & \texttt{hands \&\& !feet} & 34\,569 & 287,2 & 331,4 \\
\hline
18 & \texttt{(hotel \textbar{}\textbar{} hostel \textbar{}\textbar{} motel) \&\& (cheap \textbar{}\textbar{} budget) \&\& !luxury} & 21\,960 & 1\,468,5 & 1\,629,5 \\
\hline
19 & \texttt{(train \textbar{}\textbar{} bus \textbar{}\textbar{} ferry) \&\& (ticket \textbar{}\textbar{} pass) \&\& !airport} & 24\,660 & 2\,556,0 & 2\,679,6 \\
\hline
20 & \texttt{((castle \textbar{}\textbar{} palace) \&\& (germany \textbar{}\textbar{} austria)) \textbar{}\textbar{} (temple \&\& japan)} & 21\,299 & 1\,530,2 & 1\,661,8 \\
\hline
21 & \texttt{(the \textbar{}\textbar{} a) \&\& (of \textbar{}\textbar{} in) \&\& !zzzz} & 992\,588 & 4\,932,2 & 5\,937,3 \\
\hline
22 & \texttt{city \textbar{}\textbar{} town \textbar{}\textbar{} village \textbar{}\textbar{} hotel \textbar{}\textbar{} museum \textbar{}\textbar{} church \textbar{}\textbar{} park \textbar{}\textbar{} river \textbar{}\textbar{} beach \textbar{}\textbar{} mountain \textbar{}\textbar{} lake \textbar{}\textbar{} island \textbar{}\textbar{} road \textbar{}\textbar{} train \textbar{}\textbar{} bus \textbar{}\textbar{} airport \textbar{}\textbar{} restaurant \textbar{}\textbar{} food \textbar{}\textbar{} market \textbar{}\textbar{} street} & 790\,259 & 23\,611,6 & 24\,114,4 \\
\hline
23 & \texttt{(paris \&\& museum) \textbar{}\textbar{} louvre} & 15\,375 & 476,1 & 553,1 \\
\hline
24 & \texttt{lviv \textbar{}\textbar{} kyiv} & 3\,824 & 5,1 & 14,6 \\
\hline
\end{longtable}

}
\begin{itemize}
    \item Индекс загружается в память в 2,4 раза быстрее (840 $\to$ 348 мс; файл читается из кэша ОС, с диска разница больше) и занимает в памяти 615 МиБ вместо 1\,584.
    \item Запросы из одного-двух слов, где почти всё время --- чтение списков, стали в 2,4--3,6 раза дольше: \texttt{the} --- 278 $\to$ 679 мкс, \texttt{museum} --- 34 $\to$ 82 мкс.
    \item В сложных запросах основное время --- слияние списков и дополнение до всех документов для НЕ, а они не изменились. Поэтому разница обычно 0--20\,\% (у \texttt{!the}, где распаковывается список из 988 тыс. записей, --- 36\,\%): ИЛИ из 20 слов --- 23,6 $\to$ 24,1 мс, \texttt{!(the || and)} --- 5,2 мс у обоих. Медиана по 24 запросам выросла на 13\,\% (645 $\to$ 727 мкс), среднее --- на 7\,\%.
    \item Для утилиты командной строки, которая каждый раз загружает индекс, сжатие окупается: экономия 0,5 с на загрузке больше, чем потеря на всех 24 запросах (суммарно 4,5 мс). Для веб-сервиса, который держит индекс в памяти, запросы стали немного медленнее, зато памяти нужно в 2,6 раза меньше.
\end{itemize}

\section{Выводы}
Списки документов булева индекса сжаты кодом Variable Byte по разностям docID. На корпусе в 1 млн документов списки уменьшились в 3,3 раза (9,66 бита на запись вместо 32), весь индекс --- в 2,6 раза, с 1\,584 до 615 МиБ. Индекс загружается в 2,4 раза быстрее, проход по спискам стал в 2,6 раза дольше, а запросы --- в среднем на 7\,\% (медиана --- на 13\,\%), потому что большую часть времени запроса занимают операции над списками, а не их чтение. Корректность проверена юнит-тестами кодека и формата и сравнением со старым индексом на всём корпусе: все 364 млн записей и полная выдача 24 запросов совпали.

Недостатки и как их исправить:
\begin{itemize}
    \item \textbf{VByte не опускается ниже 8 бит на запись,} а 73\,\% записей --- у частых термов с разностями 1--10, где хватило бы 1--4 бит. Гамма- или дельта-код Элиаса, Simple-9 (несколько маленьких чисел в одном 32-битном слове) или PForDelta сжали бы эти списки в 2--4 раза сильнее; для термов из большинства документов (\texttt{the} --- в 99\,\%) выгоднее битовая карта на $N$ бит или хранение списка отсутствующих документов.
    \item \textbf{Распаковка медленнее чтения u32 в 1,7--4,2 раза,} хуже всего --- на смеси 1- и 2-байтовых кодов. Быстрее распаковываются коды без ветвления на каждый байт: Group Varint (длины четырёх чисел в одном управляющем байте), Stream VByte с SIMD.
    \item \textbf{Список распаковывается целиком,} даже если при пересечении с коротким списком нужны несколько docID. Прыжки по индексу (ЛР12) --- указатели на начало блоков сжатого списка с их первым docID --- позволят распаковывать только нужные блоки.
    \item \textbf{Пиковая память индексации не уменьшилась} (2,1 ГиБ): списки копятся в памяти массивами u32 и сжимаются только при записи. Если хранить в памяти сразу коды VByte (дописывать разность к буферу терма), память под списки сократится примерно в 3 раза, а на корпусах больше памяти нужна блочная индексация со слиянием (SPIMI), где сжатые блоки тоже меньше.
    \item \textbf{Словарь не сжат:} таблица термов (61 МиБ) и строки термов (40 МиБ) --- теперь шестая часть файла. Строки термов в виде одной строки с указателем на каждый $k$-й терм и префиксным сжатием (из лекции) заняли бы примерно вдвое меньше.
\end{itemize}
